\documentclass[10pt,a4paper]{amsart}
\usepackage[margin=19mm]{geometry}
\usepackage{amsmath,amssymb,mathtools}
\usepackage[scr=rsfs,cal=euler]{mathalfa}
\usepackage{xfrac}
\usepackage[unicode,hidelinks,bookmarks=false]{hyperref}
\newcommand{\ie}{i.e.\ }
\newcommand{\cf}{cf.\ }
\newcommand{\resp}{respectively\ }
\newcommand{\Subsection}{\subsection}
\providecommand{\nobreakdash}{\nobreak\hbox{-}}
\setlength{\emergencystretch}{2em}
\multlinegap0pt
\usepackage{bm}
\usepackage{bbm}
\usepackage{dsfont}
\usepackage{xspace}
\usepackage{cancel}
\usepackage{mathtools}
\usepackage{amsthm}
\makeatletter
\@ifundefined{theorem}{\newtheorem{theorem}{Theorem}}{}
\@ifundefined{proposition}{\newtheorem{proposition}[theorem]{Proposition}}{}
\makeatother
\title[An effective analytic recurrence for prime numbers]{An effective analytic recurrence for prime numbers}
\author{Benoit Cloitre}
\date{Source: arXiv:2508.02690v2 (CC BY 4.0)}
\begin{document}
\maketitle

\noindent\emph{Source and licence.} An effective analytic recurrence for prime numbers. Version arXiv:2508.02690v2 (CC BY 4.0), distributed under the Creative Commons Attribution 4.0 International licence, as recorded in the arXiv licence metadata for this exact version and shown on the abstract page https://arxiv.org/abs/2508.02690v2. Full rights notice: licenses/arxiv-2508.02690-CC-BY-4.0.txt.\par\smallskip

\noindent\textit{Editorial scope.} Counted unit: the Theorem whose source label is \texttt{thm:twin} in the article (source0.tex lines 401--404, source label \texttt{thm:twin}), delivered together with the article's own complete original proof of it (source0.tex lines 406--442, one \texttt{proof} environment of 37 lines). No mathematics is written, repaired or re-derived: statement and proof are the article's own text line for line. Required context retained uncounted: prop:existence\_sn (lines 216--219); thm:nagura (lines 286--289). Every definition, display and label of those lines is retained unchanged. Omitted from this package and not redistributed: 1--215, 220--285, 290--400, 405--405, 443--900. Nothing inside a retained excerpt is omitted or altered: every retained body is delivered exactly as the article's lines are written, so no omission receipt of any kind accompanies this package. Citations: the retained excerpts carry no citation command at all, so no bibliography entry of the article is transcribed and no reference is redistributed. Reference adaptation: none was needed and none was made; every reference inside the delivered unit resolves inside it, so no unresolved reference or citation is produced. Printed slips kept literal: no printed word, symbol, hypothesis or number of the counted statement or of its proof was altered. Layout and class: the article's own class is \texttt{article} with packages including fullpage, amssymb, amsmath, amsthm, color, graphicx, booktabs, float, enumitem, longtable, hyperref; the wrapper is the lane's pinned \texttt{amsart} editorial wrapper at 10pt on A4, which restates the theorem-like environments the retained text needs and adds the article's own macro definitions that the retained text needs (none), with extra wrapper packages (bm, bbm, dsfont, xspace, cancel, mathtools). The delivered numbering is the unit's own. Assets: no retained span contains a figure environment, a graphics-inclusion command, a PDF-page inclusion, a table or a TikZ picture; no image file is copied into this package, the package declares no asset dependency and no image file is redistributed with it. Licence: version arXiv:2508.02690v2 of this article is distributed under the Creative Commons Attribution 4.0 International licence, as recorded in the arXiv licence metadata for this exact version and shown on the abstract page https://arxiv.org/abs/2508.02690v2. A full rights notice is delivered at licenses/arxiv-2508.02690-CC-BY-4.0.txt. Characters: every retained excerpt is pure ASCII, so the delivered engine, XeLaTeX with its default Latin Modern text font, reports no missing character.
\begin{proposition}[Existence and characterization of $s_n^*$]
\label{prop:existence_sn}
For each $n \geq 1$, there exists a unique $s_n^* > 1$ such that $h(s_n^*) = p_{n+1} - 1$. Furthermore, $h(s) > p_{n+1} - 1$ if and only if $s > s_n^*$.
\end{proposition}

\begin{theorem}[Sharp bound]
\label{thm:nagura}
For all $n \geq 1$, we have $s_n \leq p_n$.
\end{theorem}

\begin{theorem}[Twin prime lower bound]
\label{thm:twin}
Under the twin prime conjecture, $C > \log \psi \approx 0.38225$, where $\psi$ denotes the supergolden ratio, the unique real root of $x^3 - x^2 - 1 = 0$.
\end{theorem}

\begin{proof}
By definition (Proposition \ref{prop:existence_sn}), $s_n^*$ satisfies $h(s_n^*) = p_{n+1} - 1$. Recalling $h(s) = p_{n+1}(1 + T_n^*(s))^{-1/s}$, we rearrange this equality at $s=s_n^*$ to obtain the exact equilibrium equation:
\[
1 + T_n^*(s_n^*) = \left(1 - \frac{1}{p_{n+1}}\right)^{-s_n^*}.
\]

Assuming the twin prime conjecture, consider the subsequence $\mathcal{T} = \{n : g_{n+1} = 2\}$. Define $C_{\text{twin}} := \limsup_{n \in \mathcal{T}} \sigma_n^*$. By PNT, $C_{\text{twin}} = \limsup_{n \in \mathcal{T}} s_n^*/p_{n+1}$.

We analyze the asymptotic behavior as $n \to \infty$ along a subsequence achieving the limsup.
\[
\left(1 - \frac{1}{p_{n+1}}\right)^{-s_n^*} = \exp\left(-s_n^* \log\left(1 - \frac{1}{p_{n+1}}\right)\right) = \exp\left(\frac{s_n^*}{p_{n+1}} + O\left(\frac{s_n^*}{p_{n+1}^2}\right)\right).
\]
Since $s_n^* = O(p_n)$ (Theorem \ref{thm:nagura}), the error term vanishes. Taking the limit, the equilibrium equation becomes
\[
\lim_{n\to\infty, n\in\mathcal{T}} (1 + T_n^*(s_n^*)) = \exp(C_{\text{twin}}).
\]

We analyze $T_n^*(s)$. In the limiting regime as $n \to \infty$ along the subsequence $\mathcal{T}$, composite contributions to $T_n^*(s_n^*)$ decay exponentially faster than prime contributions. The smallest composite coprime to $P_n$ exceeding $p_{n+1}$ is at least $p_{n+1}^2$. Since $s_n^* = O(p_n)$, the term $(p_{n+1}/k)^{s_n^*}$ for $k \geq p_{n+1}^2$ vanishes exponentially fast (e.g., $(1/p_{n+1})^{O(p_n)}$). Thus, only prime contributions remain in the limit.

$T_n^*(s) \geq \sum_{j=1}^{\infty} (1 + S_j/p_{n+1})^{-s}$. In the limit,
\[
T_n^*(s_n^*) \to H(C_{\text{twin}}) := \sum_{j=1}^{\infty} e^{-C_{\text{twin}} \cdot \overline{G}_j},
\]
where $\overline{G}_j$ are the limiting normalized cumulative gaps. The equilibrium is $\exp(C_{\text{twin}}) = 1 + H(C_{\text{twin}})$.

For twin primes ($g_{n+1} = 2$), admissibility modulo 3 forces $g_{n+2} \geq 4$. If $p_{n+1}, p_{n+1}+2$ are primes $>3$, then $p_{n+1} \equiv 2 \pmod 3$. Then $p_{n+1}+4 \equiv 0 \pmod 3$. Thus $S_1 = 2$ and $S_2 \geq 6$. The densest admissible pattern is $(2, 4, 2, 4, \ldots)$.

Therefore, $H(C) \geq e^{-2C} + e^{-6C} + e^{-8C} + \cdots > e^{-2C}$.

From $\exp(C) = 1 + H(C) > 1 + e^{-2C}$, we obtain $e^{3C} - e^{2C} - 1 > 0$.

Let $y = e^C > 1$. Define $f(y) = y^3 - y^2 - 1$. $f(y)$ is strictly increasing for $y > 2/3$. The unique real root is the supergolden ratio $\psi$. By Cardano's formula,
\[
\psi = \frac{1}{3}\left(1 + \sqrt[3]{\frac{29+3\sqrt{93}}{2}} + \sqrt[3]{\frac{29-3\sqrt{93}}{2}}\right) \approx 1.46557.
\]
Since $f(y) > 0$ requires $y > \psi$, we conclude $e^C > \psi$, giving $C > \log \psi \approx 0.38225$.
\end{proof}

\end{document}
